\section{Study 3: what allocation actually responds to}\label{sec:qualifier}

Study 1 contrasts a drifted memory against a clean one carrying its caveat, and
reads the difference as detection of the omission. Auditing the construction we
found that contrast varies four things at once. Appending the caveat makes the
target body 221 characters against 123--132 for every other memory --- $1.73\times$
--- and makes it the only memory in the set carrying a \textit{However} clause.
Length, surface qualification, and uniqueness move with correctness.

So we separate them. Four arms run concurrently, differing only in what is
appended to the corrupted memory's body, matched to the true caveat on character
length and on clause structure:

\begin{itemize}\itemsep1pt
\item \textbf{drifted} --- nothing appended (129 characters).
\item \textbf{clean-positive} --- a positive elaboration, no hedge (220).
\item \textbf{clean-neutral} --- a measurement-hygiene note that weakens
confidence in the finding without saying the action is risky (221).
\item \textbf{clean-negative} --- the true caveat (221).
\end{itemize}

\begin{table}[t]
\centering\small
\caption{Study 3. Verification of the corrupted memory, first pass, budget two,
by what its body carries. Bodies are matched on length and clause structure
across the three appended arms; only the content differs.}
\label{tab:qualifier}
\begin{tabular}{lrrrr}
\toprule
& drifted & clean-positive & clean-neutral & clean-negative \\
& (no qualifier) & (positive) & (weakens confidence) & (true caveat) \\
\midrule
verified & \FourDriftK/\FourDriftN{} & \FourPosK/\FourPosN{} & \FourNeutK/\FourNeutN{} & \FourNegK/\FourNegN{} \\
& (\FourDriftPct\%) & (\FourPosPct\%) & (\FourNeutPct\%) & (\FourNegPct\%) \\
\bottomrule
\end{tabular}
\end{table}

Table~\ref{tab:qualifier} orders the arms in a way the salience account does not
predict. The true caveat does not raise verification; it \emph{collapses} it
(\FourNegPct\%). A matched-length qualifier that does not answer the question
raises it (\FourNeutPct\% and \FourPosPct\%). And the silently drifted memory
sits below both (\FourDriftPct\%).

The mechanism is visible in what the budget goes to instead. Holding one model
fixed, the first-pass intent was \texttt{simplify\_onboarding} in every episode
of all four arms --- the plan does not move --- yet the allocation does. With
the true caveat present the two credits go to the memories backing that plan and
the corrupted memory is never requested. With a qualifier that raises a question
it does not settle, a credit is spent on it even though the agent does not
intend to act on it.

Two things follow.

\paragraph{The published contrast holds, but not for the stated reason.}
Verification really is higher when the caveat is absent than when it is present
($\FourNeutNegCMH$ and $\FourPosNegCMH$ for the matched arms against the true
caveat). But this is not the agent noticing that something was removed. It is
the agent not needing to look at a belief whose constraint is already written
down. Detection of the omission is not required to produce the effect, and
Study 3 gives no evidence that it occurs.

\paragraph{Silence is stealthier than qualification.} A memory that hedges about
something irrelevant draws more scrutiny than one whose material caveat was
deleted ($\FourNeutDriftCMH$). Consolidation that drops a clause therefore does
not merely remove information --- it removes the surface cue that would have
made the record look worth checking. The record becomes cleaner-looking and
correspondingly less audited, which is the worst combination for a store that
future sessions will inherit.
